\documentclass{abstract}

% Bibliography
\usepackage[
    backend=biber,
    style=authoryear,
    doi=true,
    sorting=none,
    sortcites=true,
    maxbibnames=2,
    minbibnames=1,
    maxcitenames=2,
    mincitenames=1,
    citestyle=numeric-comp,
    giveninits=true,
    isbn=false,
    date=year
]{biblatex}
\addbibresource{abstract.bib}
% For email linking
\usepackage{hyperref}

\title{ENFP429 NIST Abstract}

\author{Jacob Myers}

\keywords{Cone Calorimeter, Heat Release Rate, Composite}
\begin{document}
\maketitle
\section{Abstract}
 The overarching goal of this project is to take data from cone calorimetry tests performed by NIST/the NBS over the past 40 years, currently stored in a physical form, and create a digital, searchable database. The intent is that if someone has a material that they wish to know the data for, they can easily find tests performed on that material in the past, along with the relevant test conditions and results.

 The first step of this process is digitization. The physical data sheets are scanned, and then AI tools are used to take the data from the scanned image and create a digital data file. This file is then parsed into two smaller files. One is referred to as the data file, which contains the data resulting from the tests, like mass and energy emission, and the other is the metadata file, which contains the details of the test, like what energy input level was used, what the starting mass and time were, and the orientation of the material. The next step is to check these data and metadata files for mistakes. The process of checking the data and metadata files against the original scanned documents and editing them to remove transcription errors is termed “SmURFing.” An additional part of SmURFing is creating a Material ID, a name for the files that contains the material and primary testing conditions so that anyone can easily find the test that best matches what they are looking for.
 
 My primary duty in my role is performing this SmURFing. I access data and metadata files for a given test, and check it against the information from the original scan. I edit it as necessary, and add a Material ID. I then upload this changed version to the central NIST repository for it to be manually checked over by NIST employees. My particular data has been on composite materials. 
 
 We have additionally been educated on the function of a cone calorimeter, as well as the type of data it collects, and how this data is usually represented. We have had additional instruction on how to identify patterns of mass loss curves in cone calorimeter data. We have learned how to work with GitHub, in connecting to a central repository, pulling it down, and uploading to it. LaTeX has been used, and we have learned how to utilize LaTeX to create a document. 
 
 My SmURFing activities have been primarily focused on two technical reports. One of these reports is NBSIR 88-3733, which primarily concerns resin composites with fiber reinforcements. The resins and fibers were given within the report by their generic classifications, as provided by the material's source, which is usually the manufacturer of that specific composite. The majority of the tests used glass fibers within the composite, besides one sample of Ryton. 
 
 The two measurements from the report that were prioritized in our SmURFing were time delay delay until ignition (t\_ign) and rate of heat release (HRR). t\_ign concerned the time until pilot ignition of a panel of a given material at a given external heat flux, done by heating with a cone calorimeter in the presence of an additional hot spot. Rate of heat release was determined by analyzing oxygen consumption. This value was calculated for all of the materials multiple times at each level of external heat flux, and then graphed. These measures are meant to understand the resistance of a material to igniting in the presence of a fire or heat source and their contribution to further spread, respectively.  
 
 Other qualities were measured as well, though we did not engage with them as strongly. The heat of combustion and heat of pyrolysis were utilized to create a thermal sensitivity index, a measure of flammability. An external sensitivity index was made to indicate if a flame were self sustaining in the absence of external flux. Product (hydrocarbon, CO, and CO2) generation was also measured, and normalized to serve as a shorthand measure of combustion efficiency.
  
 The other report that I focused on was ADA168251. The goal of this report was to analyze various composite panels for their practicality of use as aircraft panel materials. The composites tested within the report were all of similar construction, being made of a woven fiber of fiberglass, Kevlar, or graphite and impregnated with either epoxy or phenolic resin and cured into a 1/8" honeycomb. They also had facings of the material making up their fibers, and a 2 mm white Tedlar covering on one side. 
 
 Similarly to the previous report, the most important data to us from this report were the ignition time (t\_ign) and Heat Release Rate. The set up for the ignition time test is very similar to the one in the previous report, with t\_ign being the measure of time until ignition in a specific orientation, at a specific external heat flux, with a exposure to a pilot flame. Testing for heat release rate was also done in a similar manner to the first report. HRR data was taken via oxygen consumption in the same manner. This procedure does explicitly have edge protection for the horizontal tests, as edge burning is not a concern in the panel's actual use case. 
 
 This report also collected data on flames on vertical panels of the relevant materials. Flame spread at the various test fluxes was measured, and compared to previously experimental data. Additionally, flame height and heat transfer of flames on vertical panels was determined for each material type. 
 
 


 Report Summaries:\\
 NBSIR 88-3733: Cone Calorimeter Evaluation of the Flammability of Composite \\
 Brown, J.E. and Braun, E. and Twilley, W.H. \\
 Expected Number of Tests: 50 \\
 Files SmURFed Successfully: 33 \\
 Files Not SmURFed Successfully: 17\\
 Tests not able to be located: 1374, 1375, 1385, 1386, 1387, 1388, 1389, 1390, 1391, 1392, 2311, 2313\\
 Files that have been re-llama-ed, but tests are not processed: 1027, 0744, 0743\\
 Files that need to be re-llama-ed: 0543
 
 
 ADA168251: The Role of Aircraft Panel Materials in Cabin Fires and their Properties \\
 Quintiere, J.G. and Babrauskas, V. and Cooper, L. and Harkleroad, M. and K. Steckler \\
 Expected Number of Tests: 68 \\
 Files SmURFed Successfully: 60 \\
 Files Not SmURFed Successfully: 8 \\
 Files that need to be re-llama-ed: 0891, 1791, 1782
 
 
 
\begin{table}[htbp]
    \centering
    \caption{NBSIR 88-3733 Combustion Test Results}
    \label{tab:combustion_results}
    \begin{tabularx}{\textwidth}{|L|c|c|c|}
        \hline
        \tableheader{Series Name} & \subunithead{t}{ign}{s} & \subunithead{t}{flameout}{s} & \unithead{Residue Yield}{\%} \\
        \hline
        Koppers6692T\_25 & 263 & NR* & NR  \\
        \hline
        Koppers6692T\_35 & 120(10)** & NR & NR \\
        \hline
        Koppers6692T\_50 & 60(2) & NR & NR \\
        \hline
        Koppers6692T\_75 & 21 & NR & NR \\
        \hline
        DowDerakanePanel\_35 & 129 & NR & NR \\
        \hline
        DowDerakanePanel\_50 & 55 & NR & NR \\
        \hline
        DowDerakanePanel\_75 & 45(2) & NR & NR \\
        \hline
        DowDerakaneHousing\_35 & 390 & NR & NR \\
        \hline
        DowDerakaneHousing\_50 & 187 & NR & NR \\
        \hline
        DowDerakaneHousing\_75 & 55(2) & NR & NR \\
        \hline
        CorflexPanel\_35 & 92(5) & NR & NR \\
        \hline
        CorflexPanel\_50 & 54(2) & NR & NR \\
        \hline
        CorflexPanel\_75 & 25 & NR & NR \\
        \hline
        Rayton1\_35 & 183 & NR & NR \\
        \hline
        Rayton1\_50 & 66 & NR & NR \\
        \hline
        Rayton1\_75 & 27(2) & NR & NR \\
        \hline
        Rayton2\_35 & 191(7) & NR & NR \\
        \hline
        Rayton2\_50 & 97 & NR & NR \\
        \hline
        Rayton2\_75 & 28(5) & NR & NR \\
        \hline
        Rayton3\_25 & 154(7) & NR & NR \\
        \hline
        Rayton3\_50 & 75 & NR & NR \\
        \hline
        Rayton3\_75 & 29(1) & NR & NR \\
        \hline
        Rayton4\_35 & N.I.*** & NR & NR \\
        \hline
        Rayton4\_50 & 88 & NR & NR \\
        \hline
        Rayton4\_75 & 38(3) & NR & NR \\
        \hline
    \end{tabularx}
\end{table}

\begin{table}[htbp]
	\centering
	\caption{ADA168251 Combustion Test Results}
	\label{tab:combustion_results_2}
	\begin{tabularx}{\textwidth}{|L|c|c|c|}
		\hline
		\tableheader{Series Name} & \subunithead{t}{ign}{s} & \subunithead{t}{flameout}{s} & \unithead{Residue Yield}{\%} \\
		\hline
		EpoxyFiberglass\_Hor\_25 & 31.6 & NR & NR \\
		\hline
		EpoxyFiberglass\_Hor\_50 & 7.8 & NR & NR \\
		\hline
		EpoxyFiberglass\_Hor\_75 & 4.8 & NR & NR \\
		\hline
		EpoxyFiberglass\_Vert\_25 & 29.8 & NR & NR \\
		\hline
		EpoxyFiberglass\_Vert\_50 & 7.9 & NR & NR \\
		\hline
		EpoxyFiberglass\_Vert\_75 & 5.8 & NR & NR \\
		\hline
		PhenolicFiberglass\_Hor\_25 & 27.9 & NR & NR \\
		\hline
		PhenolicFiberglass\_Hor\_50 & 8.0 & NR & NR \\
		\hline
		PhenolicFiberglass\_Hor\_75 & 4.3 & NR & NR \\
		\hline
		PhenolicFiberglass\_Vert\_25 & N.I. & NR & NR \\
		\hline
		PhenolicFiberglass\_Vert\_50 & 8.0 & NR & NR \\
		\hline
		PhenolicFiberglass\_Vert\_75 & 5.5 & NR & NR \\
		\hline
		EpoxyKevlar\_Hor\_25 & 33.0 & NR & NR \\
		\hline
		EpoxyKevlar\_Hor\_50 & 8.8 & NR & NR \\
		\hline
		EpoxyKevlar\_Hor\_75 & 3.9 & NR & NR \\
		\hline
		EpoxyKevlar\_Vert\_25 & 36.1 & NR & NR \\
		\hline
		EpoxyKevlar\_Vert\_50 & 6.5 & NR & NR \\
		\hline
		EpoxyKevlar\_Vert\_75 & 5.5 & NR & NR \\
		\hline
		PhenolicKevlar\_Hor\_25 & N.I. & NR & NR \\
		\hline
		PhenolicKevlar\_Hor\_50 & 8.0 & NR & NR \\
		\hline
		PhenolicKevlar\_Hor\_75 & 4.5 & NR & NR \\
		\hline
		PhenolicKevlar\_Vert\_25 & N.I & NR & NR \\
		\hline
		PhenolicKevlar\_Vert\_50 & 8.5 & NR & NR \\
		\hline
		PhenolicKevlar\_Vert\_75 & 5.5 & NR & NR \\
		\hline
		PhenolicGraphite\_Hor\_25 & N.I. & NR & NR \\
		\hline
		PhenolicGraphite\_Hor\_50 & 12.2 & NR & NR \\
		\hline
		PhenolicGraphite\_Hor\_75 & 5.5 & NR & NR \\
		\hline
		PhenolicGraphite\_Vert\_25 & N.I. & NR & NR \\
		\hline
		PhenolicGraphite\_Vert\_50 & 9.5 & NR & NR \\
		\hline
		PhenolicGraphite\_Vert\_75 & 5.5 & NR & NR \\
		\hline
	\end{tabularx}
\end{table}
 
 *NR is used when data was unavailable within the report \\
 **The number in parentheses is the variance when multiple t\_igns were measured \\
 ***N.I. is used when no ignition was observed \\
 
 \cite{ADA168251, NBSIR883733}




\section{References}
\printbibliography[heading=none]

\end{document}